% Research contribution of agent /root/monodromy.  Input into a parent document.
% Status is local to the assertions below; this is not a claim to solve all of A.3--A.5.
\section{A.3--A.5: proved statements and remaining gates}
\label{sec:monodromy-contribution}
\status{Incomplete universal determination; the individually labelled results below are proved.  Lead A's full-linear assertion is refuted, its positivity assertion is proved with an explicit threshold, and its plane-curve point-count consequence is proved.}

\subsection{Conventions and the generic rank for projective hypersurfaces}

Let $S_d=\mathbf P(H^0(\mathbf P^r,\mathcal O(d)))$, and let
$\pi:\mathcal Y_d\to S_d$ be the universal degree-$d$ hypersurface.
For $r\geq2$, the sheaf $\mathcal H=R^{r-1}\pi_*\mathcal O_{\mathcal Y_d}$
is locally free of rank
\[
 h_d=\begin{cases}\binom{d-1}{r}&d\geq r+1,\\0&d\leq r.\end{cases}
\]
This follows from the universal divisor exact sequence and the cohomology
of line bundles on projective space; formation of $\mathcal H$ commutes
with all base changes.  Absolute Frobenius induces a map
$F_S^*\mathcal H\to\mathcal H$.  On its invertibility locus $W_d$, the
kernel of $F-1$ is a locally constant $\mathbf F_p$-sheaf of rank $h_d$.
On the smooth locus it is $R^{r-1}\pi_*\mathbf F_p$, by the
Artin--Schreier sequence.  For $X=\mathbf P^r$, this is the vanishing
local system in the question, since $H^{r-1}(\mathbf P^r,\mathcal O)=0$.

\begin{proposition}[All degrees, all characteristics]\label{prop:projective-generic-hw-rank}
The generic degree-$d$ hypersurface of $\mathbf P^r$ has bijective
Frobenius on $H^{r-1}(\mathcal O)$, and its generic stable rank is $h_d$.
In particular, for plane curves the generic rank is
$g_d=(d-1)(d-2)/2$, and the smallest positivity threshold is $d_0=3$.
\end{proposition}

\begin{proof}
Work over an algebraic closure and take $d$ hyperplanes in general
position.  Their union $D$ is a simple normal crossings divisor.
The alternating resolution of $\mathcal O_D$ by the structure sheaves
of intersections has only $H^0$ in its first hypercohomology page,
because every nonempty intersection is a projective space.  Frobenius
is bijective on each of these $H^0$'s and hence on the cohomology of
their complex.  It is therefore bijective on
$H^{r-1}(D,\mathcal O_D)$.  Thus $W_d$ is nonempty.  The smooth locus
is a nonempty dense open of the irreducible scheme $S_d$, so it meets
$W_d$.  Invertibility is open, which proves the assertion.  This proof
establishes only bijectivity on the specified coherent cohomology;
it does not conflate this condition with every formulation of
ordinarity for higher dimensional varieties.
\end{proof}

\paragraph{Comparison with Illusie.}
Illusie (1990), Theorem~0.1, p.~375, proves the stronger statement that
the general smooth complete intersection of every specified dimension
and multidegree is ordinary: its Newton and Hodge polygons coincide.
The theorem and the explanation of the earlier Koblitz argument on
pp.~375--376 were inspected in the publisher's two-page preview of
the reprint.  In particular, generic Hasse--Witt bijectivity follows;
the proposition above gives a separate complete proof of exactly the
part needed here.  The parameter space in Illusie's theorem varies all
defining equations.  It does not, by restriction of an open set alone,
prove generic ordinarity for hypersurface sections of every fixed $X$.

\subsection{The determinant character and the Hasse invariant}
\label{subsec:hasse-determinant}

Let $S$ be connected and let $(\mathcal E,F)$ be a vector bundle of
rank $h$ in characteristic $p$, with $F:F_S^*\mathcal E\simeq\mathcal E$.
Write $V=\ker(F-1)$ and $\mathcal L=\det\mathcal E$.  Taking
determinants gives an invertible section
\[
 \operatorname{Ha}_F\in
 H^0\bigl(S,\operatorname{Hom}(F_S^*\mathcal L,\mathcal L)\bigr)
 =H^0(S,\mathcal L^{\otimes(1-p)}).
\]
This is the Hasse invariant of the specified Frobenius module.  It
determines the determinant character exactly, including its arithmetic
part: $\det V$ is the rank-one Lang local system
$\ker(F_{\mathcal L}-1)$.  If $e$ trivializes $\mathcal L$ and
$F_{\mathcal L}(e)=h e$, a fixed nonzero vector $ze$ satisfies
\begin{equation}\label{eq:determinant-kummer}
 z^{p-1}=h^{-1}.
\end{equation}
Thus the $\mathbf F_p^\times$-torsor of frames of $\det V$ is the
Kummer torsor \eqref{eq:determinant-kummer}; a deck transformation
$z\mapsto az$ acts on the determinant fibre by $a$.  Changing the
trivialization of $\mathcal L$ changes this equation by a
$(p-1)$st power, as required.  Over $\mathbf F_q$, geometric Frobenius
acts on $\det V$ by the determinant of the $q$-power coherent
Frobenius; arithmetic Frobenius is its inverse.  This description
holds for the stable bundle of the relative vanishing system on its
constant-rank locus as well as for the full ordinary coherent module.
It does not assert that the image of the Kummer character is always
all of $\mathbf F_p^\times$: that depends on its Kummer class.

\subsection{The exact plane-curve monodromy in degrees at most four}
\label{subsec:small-plane-monodromy}

For $d=1,2$, the rank is zero and both monodromy groups are trivial.
For $d=3$, both images are $\mathrm{GL}_1(\mathbf F_p)$: for odd $p$
this follows from the Legendre subfamily and its simple Hasse roots,
proved below; for $p=2$ the assertion means the trivial group on a
one-dimensional space.  A nonzero rank-one representation is
irreducible.

For $d=4$, both images are $\mathrm{GL}_3(\mathbf F_p)$, and the
representation is irreducible.  Here is why the moduli-space theorem
actually applies.  Every nonhyperelliptic smooth genus-three curve is
canonically a smooth plane quartic, and conversely.  The stack quotient
of the parameter space of smooth plane quartics by $\mathrm{PGL}_3$
is the nonhyperelliptic open substack of $\mathcal M_3$.
The map from the parameter space to this quotient is a
$\mathrm{PGL}_3$-torsor in the stack sense, with geometrically connected
fibres, and is surjective on etale fundamental groups: a connected
finite etale cover of the quotient remains connected after pullback
along a torsor under a connected algebraic group.  Restriction to
the ordinary locus preserves this assertion.  That locus is a
nonempty open of the ordinary locus of $\mathcal M_3$, hence its
fundamental group has the same finite monodromy image.
Achter--Pries (2008), Proposition~4.8, says that this image on the
etale $p$-torsion of the Jacobian is $\mathrm{GL}_3(\mathbf F_p)$.
Dualizing gives the same full group on $V$.
The arithmetic image contains the geometric image and is already
a subgroup of the same finite general linear group, so it too is full.
No analogous moduli-space surjectivity assertion is used for degrees
$d\geq5$.  In characteristic two the separate contact-curve argument
and group induction instead determine every degree, in
Theorem~\ref{thm:plane-char2-all-groups}.

\subsection{A universal obstruction in characteristic two}
\label{subsec:odd-plane-obstruction}

Let $p=2$ and let $d\geq3$ be odd.  Write
\[
 f=\sum_{|a|=d}c_aX^{a_1}Y^{a_2}Z^{a_3},\qquad
 I_d=\{u\in\mathbf Z_{>0}^3:|u|=d\}.
\]
Put $c_a=0$ if a coordinate of $a$ is negative.  On a coefficient
chart of $S_d$, use the Cech basis of
$H^2(\mathbf P^2,\mathcal O(-d))$ indexed by $I_d$.
The matrix of the semilinear Frobenius on $H^1(C_f,\mathcal O)$ is
\begin{equation}\label{eq:plane-hw-char2}
 H_{u,v}=c_{2v-u}\quad(u,v\in I_d),\qquad F(x)=Hx^{[2]}.
\end{equation}
Indeed, the connecting map from
$0\to\mathcal O(-d)\xrightarrow{f}\mathcal O\to\mathcal O_{C_f}\to0$
identifies Frobenius with multiplication by $f$ after squaring the
Cech representative $X^{-v_1}Y^{-v_2}Z^{-v_3}$.

Define a column $w=(w_u)$ by
\[
 w_u=\begin{cases}c_u&u_1,u_2,u_3\text{ all odd},\\0&\text{otherwise}.
 \end{cases}
\]
Then the following is an identity over the polynomial ring in the
coefficients, not merely a numerical identity at selected curves:
\begin{equation}\label{eq:universal-covector}
 H^t w=w^{[2]}.
\end{equation}
For its $v$-coordinate, pair the term $c_{2v-u}c_u$ indexed by $u$
with the term indexed by $2v-u$.  A nonzero term has both indices
positive and all odd, so the involution preserves the indexing set.
All two-element orbits cancel in characteristic two.  The only possible
fixed index is $u=v$, contributing $c_v^2$ exactly when $v$ is all odd.

If $H$ is invertible, then $w\ne0$.  Otherwise all coefficients with
three odd exponents vanish, and every row with all-odd index $u$
in \eqref{eq:plane-hw-char2} vanishes.  Such an index exists because
$(1,1,d-2)\in I_d$.  This contradicts invertibility.

For a Lang fixed vector $x$, namely $Hx^{[2]}=x$, equation
\eqref{eq:universal-covector} gives
\[
 (w^tx)^2=(w^{[2]})^tx^{[2]}=w^tHx^{[2]}=w^tx.
\]
Thus $x\mapsto w^tx$ is an $\mathbf F_2$-valued linear functional on
the local system $V=\ker(F-1)$.  It is nonzero, since the Lang fixed
vectors form a basis after an etale cover.  It is a section of
$V^\vee$, so it is fixed by both geometric and arithmetic monodromy.
Working on a dense coefficient chart suffices: its intersection with
the ordinary smooth locus is dense, and the map on fundamental groups
from a nonempty open of a normal integral scheme is surjective.

\begin{theorem}[Refutation of full linear monodromy]\label{thm:odd-plane-char2-obstruction}
For every finite field of characteristic two and every odd $d\geq5$,
let $\bar\eta$ be a geometric generic point of the ordinary smooth
degree-$d$ plane-curve locus $U_d$, and put
$V=H^1_{\mathrm{et}}(C_{\bar\eta},\mathbf F_2)$.
Both images
\[
 G_{\mathrm{geom}}=\mathrm{im}\bigl(\pi_1(U_d\otimes\bar k,\bar\eta)
 \to\mathrm{GL}(V)\bigr),\qquad
 G_{\mathrm{arith}}=\mathrm{im}\bigl(\pi_1(U_d,\bar\eta)
 \to\mathrm{GL}(V)\bigr)
\]
fix a nonzero covector.  Consequently both are contained in the
proper subgroup
\[
 \mathrm{Stab}(\ell)\simeq
 \mathbf F_2^{\,g_d-1}\rtimes\mathrm{GL}_{g_d-1}(\mathbf F_2),
\]
and $V$ is reducible, with invariant hyperplane $\ker\ell$.
There is no integer $d_1(\mathbf P^2)$ for which the proposed full
linear geometric or arithmetic monodromy holds for every $d\geq d_1$
in characteristic two.
\end{theorem}

The theorem proves containment, not equality with the displayed
stabilizer.  Determining the exact image remains a separate gate.
The determinant character in this characteristic is automatically
trivial, because $\mathbf F_2^\times=\{1\}$.

\paragraph{Geometric interpretation and literature check.}
The obstruction agrees with the difference between the effective
geometric theta characteristic
$\mathcal O_C((d-3)/2)$ and the canonical characteristic-two theta
characteristic, which is non-effective for an ordinary curve.
Cais, Ellenberg, and Zureick-Brown (2013, Theorem~4.2) prove the resulting
rational two-torsion statement for odd degree plane curves.
Castryck, Streng, and Testa (2014, Corollary~13, Remark~14, and
Theorem~16) give the ordinary and Hasse--Witt formulations.
The explicit identity above avoids any descent ambiguity for theta
characteristics over imperfect coefficient fields.

\paragraph{Explicit checked fibre.}
Over $\mathbf F_2$, take
\begin{align*}
f={}&x^5+x^4y+x^3y^2+x^2y^3+x^2y^2z+xy^4+xyz^3+xz^4\\
 &+y^4z+y^3z^2+yz^4.
\end{align*}
In the ordered basis
$(1,1,3),(1,2,2),(1,3,1),(2,1,2),(2,2,1),(3,1,1)$, its matrix is
\[
H=\begin{pmatrix}
1&0&0&0&0&0\\1&0&1&0&1&1\\0&1&0&0&0&0\\
1&1&0&0&1&1\\0&0&1&0&1&0\\0&0&0&1&0&0
\end{pmatrix},\qquad w=(1,0,0,0,0,0)^t.
\]
Its rank is six.  The ideal $(f,f_x,f_y,f_z)$ has Groebner basis
$\{1\}$ on each of the three standard projective charts, proving
geometric smoothness.  Direct enumeration over
$\mathbf F_{2^m}$ for $m=1,\ldots,6$ gives respectively
$3,3,9,23,38,33$ points; the traces of $H^m$ modulo two are
$0,0,0,0,1,0$.  The program \texttt{code/plane\_char2.py} and its JSON
output record the input, field arithmetic, certificates, and timings.
The fibre is a check of the formulas; the universal identity, rather
than this one fibre, proves the generic monodromy obstruction.

\subsection{Plane curves: every point-count residue in every open}
\label{subsec:plane-residues}

\begin{theorem}\label{thm:plane-all-residues}
Let $k$ be a finite field of characteristic $p$.  For every $d\geq4$,
every nonempty Zariski open $U\subset S_d$, and every $a\in\mathbf F_p$,
there exist a finite extension $k'/k$ and a smooth plane curve
$C\in U(k')$ of degree $d$ with $\#C(k')\equiv a\pmod p$.
For $d=1,2$ the only residue for smooth curves is $1$.
For $d=3$, the residues on the ordinary smooth locus are exactly
$\mathbf F_p\setminus\{1\}$, and every nonempty open of that locus
attains every such residue after a finite extension.  Supersingular
smooth cubics give residue $1$, so all residues occur among all smooth
cubics; the every-open assertion for all residues fails in degree three.
\end{theorem}
\begin{proof}

\paragraph{Elliptic input.}
For every $\lambda\in\mathbf F_p^\times$, there is an ordinary elliptic
curve $E/\mathbf F_Q$, with $\mathbf F_Q/k$ finite, for which geometric
Frobenius on $H^1_{\mathrm{et}}(E_{\bar k},\mathbf F_p)$ is multiplication
by $\lambda$.  For $p=2$ take any ordinary elliptic curve.  For odd $p$,
write $m=(p-1)/2$.  The Legendre family
$y^2=x(x-1)(x-t)$ has Hasse polynomial
\[
 A(t)=(-1)^m\sum_{i=0}^{m}\binom mi^2t^i.
\]
Its roots are simple and avoid $0,1$: its values there are nonzero,
and it satisfies
$t(1-t)A''+(1-2t)A'-m^2A=0$ modulo $p$.
A repeated root away from $0,1$ would, by successive differentiation,
force all derivatives through degree $m<p$ to vanish, a contradiction.
At a simple root, the Kummer cover $z^{p-1}=A(t)^{-1}$ is totally
ramified of degree $p-1$.  Its geometric monodromy is therefore the
whole $\mathbf F_p^\times$.  Finite-field Chebotarev on the ordinary
Legendre locus gives every scalar $\lambda$.
Moreover we may replace $Q$ by $Q^e$ with arbitrarily large
$e\equiv1\pmod{p-1}$ without changing $\lambda$.

\paragraph{A nodal curve with prescribed residue.}
Put $r=d-3\geq1$ and embed $E$ as a plane cubic by $|3O|$.
Choose rational lines $L_1,\ldots,L_r$ in general position such that
all intersections are transverse, no three components meet, the
intersection $E\cap L_i$ is split for $i\geq2$, and
$E\cap L_1$ has either cycle type $3$ or cycle type $1+2$ under
Frobenius.  Write $t=0$ or $1$ for the number of rational points in
this last intersection.  Such lines exist after an extension of the
above kind, as follows.

Split lines arise from distinct $P,Q,R\in E(\mathbf F_Q)$ with
$P+Q+R=0$, giving order $Q^2$ choices.  Lines of type $1+2$ arise from
$P\in E(\mathbf F_{Q^2})\setminus E(\mathbf F_Q)$ and
$R=-(P+\pi P)\in E(\mathbf F_Q)$; the three points are distinct,
collinear, and Frobenius invariant.  There are order $Q^2$ choices.
For type $3$, use nonrational points in
$\ker(1+\pi+\pi^2)$.  This is the kernel of a separable isogeny,
is contained in $E(\mathbf F_{Q^3})$, and has cardinality
$\#E(\mathbf F_{Q^3})/\#E(\mathbf F_Q)$, of order $Q^2$.
At most nine of these points are rational, since rational points in
the kernel lie in $E[3]$.  Thus there are order $Q^2$ suitable lines
of each required type.  The finitely many incidence and tangency
conditions excluded at each choice remove only order $Q$ rational
lines.  This proves existence for sufficiently large $Q$.

The connected nodal curve
$C_0=E\cup L_1\cup\cdots\cup L_r$ has degree $d$.
Its normalization sequence
\[
0\longrightarrow\mathcal O_{C_0}\longrightarrow
\mathcal O_E\oplus\bigoplus_i\mathcal O_{L_i}
\longrightarrow\bigoplus_{P\text{ node}}k(P)\longrightarrow0
\]
shows that Frobenius on $H^1(C_0,\mathcal O_{C_0})$ is bijective:
it is bijective on the finite reduced node scheme, on constants of
the normalization, and on $H^1(E,\mathcal O_E)$.
Inclusion--exclusion gives
\begin{align*}
 \#C_0(\mathbf F_Q)
 &\equiv (1-\lambda)+r-\{t+3(r-1)\}-\binom r2\\
 &\equiv 4-2r-\binom r2-t-\lambda\pmod p.
\end{align*}
As $t\in\{0,1\}$ and $\lambda\in\mathbf F_p^\times$ vary, these
residues fill $\mathbf F_p$, including when $p=2$.

\paragraph{Passage to an arbitrary smooth open.}
The ordinary coherent-cohomology locus $W_d\subset S_d$ introduced
above is geometrically integral and contains $C_0$.
The finite etale local system $\ker(F-1)$ is defined over all of $W_d$,
including its singular fibres.  After base changing to $\mathbf F_Q$,
the Frobenius class of $C_0$ belongs to the finite arithmetic
monodromy group.  Chebotarev gives that same class at a closed point
of any nonempty open of $W_d$.  Apply this to
$U\cap W_d\cap S_d^{\mathrm{smooth}}$.
The coherent Frobenius trace formula
\[
 \#C(\mathbf F_{Q^e})\equiv
 1-\operatorname{Tr}\bigl(\operatorname{Frob}_{Q^e}^{\mathrm{geom}},
 H^1_{\mathrm{et}}(C_{\bar k},\mathbf F_p)\bigr)\pmod p
\]
also applies to the proper nodal curve $C_0$, so equality of the
classes gives the desired residue.  More precisely, the class at
$C_0$ maps to degree one in the finite constant-field quotient, and
Chebotarev supplies it at degrees congruent to one in that quotient;
the permitted residue fields are finite extensions of $k$.
This argument never assumes full general linear monodromy for
degree-$d$ plane curves.
\end{proof}

\paragraph{Primary-source verification of the two inputs.}
Katz--Sarnak (1999), Theorem~9.7.13, p.~290, is the finite-monodromy
equidistribution theorem, with the constant-field coset condition made
explicit in 9.7.1--9.7.9.  To apply its characteristic-zero formulation
to our finite $\mathbf F_p$-representation, take a faithful finite
dimensional characteristic-zero representation of its finite image
(for example the regular representation), choose $\ell\ne p$, and
regard it as a lisse $\overline{\mathbf Q}_\ell$-sheaf.  It is pure of
weight zero, and its finite geometric and arithmetic groups have the
required normality.  The theorem gives a positive limiting proportion
for every nonempty conjugacy class in the required coset.  Restriction
to a nonempty open preserves the finite monodromy image because the
base is normal.  Katz (1973), SGA~7~II, Expose~XXII, Corollary~3.2,
formulas (3.2.1)--(3.2.2), pp.~410--411, gives precisely the trace
congruence for every proper scheme over a finite field, including
our nodal curve.  Proposition~2.0 and formula (2.2.4) of the same
expose identify its coherent and etale formulations.

\subsection{An explicit positivity threshold for every smooth ambient variety}
\label{subsec:explicit-positivity}

Here $X\subset\mathbf P^N$ is smooth projective of pure dimension
$n\geq2$ over a perfect field of characteristic $p$, and $L=\mathcal O_X(1)$.
Geometric connectedness is not needed for this proposition.  All
dimensions in the following definitions may be computed after extending
the ground field to an algebraic closure.  Put
\[
 D=\deg_L X,\qquad
 a=\dim H^{n-1}(X,\mathcal O_X)_{\mathrm{ss}},
\]
and, for $1\leq m\leq n-2$, define the explicit cohomological bounds
\[
 B_m=\sum_{j=0}^{m}\binom mj
       h^{n-m-1+j}(X,L^{-j}),\qquad
 B=D+a+\sum_{m=1}^{n-2}B_m.
\]
Let $r_\omega$ be any Castelnuovo--Mumford regularity bound for
$i_*\omega_X$ on $\mathbf P^N$, and put
\begin{equation}\label{eq:explicit-d0}
 d_0=\max\left\{2n,1,r_\omega-1,
       \left\lfloor\frac{2^n n!B}{D}\right\rfloor+1\right\}.
\end{equation}
Every term is determined by the embedding and finitely many coherent
cohomology groups of $X$; this is a computable bound, not an unspecified
``sufficiently large'' constant.

\begin{theorem}\label{thm:all-X-explicit-positive-rank}
For every $d\geq d_0$, the generic smooth hypersurface section $Y$ of
degree $d$ has
$\dim E(Y,\mathcal O_Y)_{\mathrm{ss}}>0$.
More precisely, for every $d\geq\max\{1,r_\omega-1\}$, its generic
rank $f_{X,d}$ satisfies
\begin{equation}\label{eq:generic-rank-lower-bound}
 f_{X,d}\ \geq\
 D\binom dn-D\binom d{n-1}
 -\sum_{m=1}^{n-2}B_m\binom dm-a,
\end{equation}
where a binomial coefficient with lower index greater than $d$ is zero.
Consequently
\[
 f_{X,d}=\frac{D}{n!}d^n+O_X(d^{n-1})\quad(d\longrightarrow\infty).
\]
\end{theorem}

\begin{proof}
Take $d$ hyperplanes in simultaneous general position over the
algebraic closure and set $Y_0=\bigcup_{i=1}^d(X\cap H_i)$.
Every intersection $Z_I=X\cap\bigcap_{i\in I}H_i$ with $|I|\leq n$
is smooth of dimension $n-|I|$, and intersections of $n+1$ components
are empty.  Intersections with $|I|=n$ consist of $D$ reduced points.
The alternating normal-crossings resolution gives the Frobenius
equivariant spectral sequence
\[
 E_1^{s,t}=\bigoplus_{|I|=s+1}H^t(Z_I,\mathcal O_{Z_I})
 \ \Longrightarrow\ H^{s+t}(Y_0,\mathcal O_{Y_0}).
\]
The term $E_1^{n-1,0}$ has dimension $D\binom dn$ and bijective
Frobenius.  It has no outgoing differential on any page.  The image
of the incoming $d_1$ has dimension at most
$D\binom d{n-1}$: each smooth positive dimensional intersection has
at most $D$ connected components, since each has positive integral
degree.  For $r=2,\ldots,n-1$, the possible incoming $d_r$ has source
$E_r^{n-1-r,r-1}$.  Writing $m=n-r$, its dimension is at most
$\binom dm B_m$.  Indeed, the Koszul resolution of a transverse
$m$-fold hyperplane section bounds
\[
 h^{n-m-1}(Z_I,\mathcal O_{Z_I})
 \leq\sum_{j=0}^{m}\binom mj
 h^{n-m-1+j}(X,L^{-j})=B_m.
\]
Every quotient of the top term by an incoming image still has
bijective Frobenius: a Frobenius-stable subspace of a finite dimensional
bijective semilinear space over a perfect field itself has bijective
Frobenius.  Hence $E_\infty^{n-1,0}$ is a bijective-Frobenius subquotient
of $H^{n-1}(Y_0,\mathcal O_{Y_0})$ of dimension at least the first three
terms of the right side of \eqref{eq:generic-rank-lower-bound}.
Exactness of the stable-part functor gives the same lower bound for
$\dim H^{n-1}(Y_0,\mathcal O_{Y_0})_{\mathrm{ss}}$.

For $d\geq\max\{1,r_\omega-1\}$, regularity and Serre duality give
$H^{n-1}(X,L^{-d})=H^1(X,\omega_X\otimes L^d)^\vee=0$.
The divisor sequence, including for the normal-crossings divisor,
then makes $h^{n-1}(Y,\mathcal O_Y)$ independent of $Y$ in the space of
properly intersecting degree-$d$ sections: the terminal map
$H^n(X,L^{-d})\to H^n(X,\mathcal O_X)$ is surjective because
$H^n(Y,\mathcal O_Y)=0$.  The universal family is
flat, and cohomology and base change therefore gives a vector bundle
with Frobenius.  If its rank is $h$, the stable rank at a perfect
geometric point is the rank of $F^h$; thus its maximal value is attained
on an open, and the generic value is at least the value at $Y_0$.
The smooth locus is dense and meets this open.

For a smooth $Y$, R1 identifies the stable ambient image with
$H^{n-1}(X,\mathcal O_X)_{\mathrm{ss}}$, of dimension $a$; in fact the
coherent restriction is injective in the present range of $d$.
Exactness of the stable-part functor subtracts $a$ and proves
\eqref{eq:generic-rank-lower-bound}.
For $d\geq2n$ one has
$\binom dn\geq(d/2)^n/n!$ and $\binom dm\leq d^{n-1}$ for $m<n$.
The right side of \eqref{eq:generic-rank-lower-bound} is therefore at
least $d^{n-1}(Dd/(2^nn!)-B)$, which is positive under
\eqref{eq:explicit-d0}.  Finally, Serre duality, Riemann--Roch, and
Serre vanishing give
$\dim E(Y,\mathcal O_Y)=Dd^n/n!+O_X(d^{n-1})$; together with the lower
bound this proves the asymptotic formula.
\end{proof}

\begin{corollary}[The eventual parenthesis of 8.5.8]
\label{cor:eventual-hodge-parenthesis}
For the explicit $d_0$ of \eqref{eq:explicit-d0}, every $d\geq d_0$
has the following properties.  Every smooth degree-$d$ section $Y$
satisfies
\[
 \dim E(Y,\mathcal O_Y)
 =h^0(X,\omega_X\otimes L^d)-h^0(X,\omega_X)>0.
\]
For every $Y$ in the maximal-stable-rank open $U_d$, the number of
$p$-adic unit roots in vanishing cohomology is positive.  Thus, on
this open and in this stated eventual range, unit-root nonvanishing
and the maximal Hodge-level condition $E(Y,\mathcal O_Y)\ne0$ are
equivalent.  No assumption on crystalline $p$-torsion of $X$ is needed.
\end{corollary}

\begin{proof}
The coherent divisor sequence and the vanishings used in the preceding
proof identify $E(Y,\mathcal O_Y)$ with
$\ker(H^n(X,L^{-d})\to H^n(X,\mathcal O_X))$, whose dimension is the
displayed, $Y$-independent difference by Serre duality.  It is positive
because it is at least the positive generic stable rank of
Theorem~\ref{thm:all-X-explicit-positive-rank}.
For $Y\in U_d$, Theorems~\ref{thm:A-exact85} and
\ref{thm:A-unit-count} identify the unit-root count with
\[
 f_{X,d}+\dim H^n_{\mathrm{cris}}(X/W)[p]_{\mathrm{ss}},
\]
which is positive.  Both properties therefore hold for every point
under these quantifiers.
\end{proof}

\paragraph{Scope of the manuscript's parenthesis.}
The supplied transcription of cote~4, scan~8, first specifies a
section for which the semisimple rank has its maximum value.  Scan~9,
continuing 8.5.8, then says ``pour tous ces $Y$'' and, equivalently,
``pour $Y$ g\'en\'erique''.  Corollary~\ref{cor:eventual-hodge-parenthesis}
answers this eventual generic assertion, including varieties with
$p$-torsion.  It does not assert that an arbitrary smooth section
outside $U_d$ has a unit root whenever its coherent vanishing space
is nonzero.  For example, a supersingular smooth plane cubic has
one-dimensional $H^1(\mathcal O)$ and no unit root.  That example
distinguishes the stronger pointwise assertion from the quantifiers
actually written in 8.5.8; it is not substituted for the requested
eventual assertion.

\paragraph{Relation to Koblitz.}
Koblitz (1975), I.\S5, Theorems~1--2, pp.~139--141, proves the
asymptotic invertibility estimate with an error of order $d^{n-1}$.
The argument above gives explicit constants in terms of the embedding
and the displayed cohomology groups.  Koblitz's I.\S6, Theorem~3,
p.~142, concerns generic complete intersections, not every hypersurface
section of an arbitrary fixed ambient variety.  His revised conjecture
in II.\S3, pp.~149--151, must not be cited as an established exact rank
formula.

\subsection{All residues for hypersurfaces in every projective space}

\begin{theorem}\label{thm:projective-all-residues}
Let $k$ be finite of characteristic $p$, let $n\geq2$, and let
$d\geq pn+1$.  Every nonempty open in the degree-$d$ hypersurface
parameter space of $\mathbf P^n_k$ contains, over a finite extension,
a smooth hypersurface with any prescribed point-count residue
modulo $p$.
\end{theorem}

\begin{proof}
Choose the integer representative $t\in\{0,\ldots,p-1\}$ of the
desired residue.  Over a sufficiently large finite extension
$\mathbf F_Q/k$, choose a general-position collection of $d$
geometric hyperplanes on which Frobenius has $t$ orbits of length $n$
and one orbit of length $d-nt>n$.  Existence follows by taking a
rational point in the open set of general-position configurations
in the product of Weil restrictions of the dual projective space
from the indicated extension degrees.  This parameter space is
geometrically integral, since over the algebraic closure it is a
product of projective spaces; the general-position open is nonempty.
Lang--Weil supplies rational points over arbitrarily large fields.
One may first form these Weil restrictions over a fixed finite base
field and then take extension degrees congruent to one modulo the
least common multiple of the prescribed orbit lengths; this preserves
their field factors while making the field arbitrarily large.
Distinctness of every geometric hyperplane forces the specified
orbit lengths to be exact.

The union $Y_0$ of these hyperplanes is defined over $\mathbf F_Q$.
A rational point lying on one hyperplane lies on its entire Frobenius
orbit.  An orbit of length $n$ has exactly one common intersection
point in $\mathbf P^n$, and that point is rational.  An orbit of
length greater than $n$ has empty intersection.  The $t$ points from
the shorter orbits are distinct, because any $n+1$ distinct
hyperplanes have empty intersection.  Therefore
\[
 \#Y_0(\mathbf F_Q)=t
\]
as an equality of integers.
The general-position hyperplane-union argument proves that
$H^{n-1}(Y_0,\mathcal O_{Y_0})$ has bijective Frobenius, of the full
generic rank $\binom{d-1}{n}$.  Thus $Y_0$ belongs to the same open
$W_d$ on which the entire Lang local system is defined.
The coherent trace congruence, together with the standard vanishing
of intermediate coherent cohomology for a projective hypersurface,
reads
\[
 \#Y(\mathbf F_{Q^e})\equiv
 1+(-1)^{n-1}\operatorname{Tr}(\operatorname{Frob}_{Q^e}^{\mathrm{geom}},V_Y)
 \pmod p.
\]
Apply finite Chebotarev on $W_d$ to the Frobenius class of $Y_0$ and
the given smooth open.  The resulting smooth hypersurface has the
same residue $t$.  No exact monodromy-group computation is required.
\end{proof}

\subsection{Relation to the subsequent classifications}
The odd-degree characteristic-two obstruction refutes the universal
full-$\mathrm{GL}$ expectation, even for geometrically connected
$X=\mathbf P^2$.  The all-degree characteristic-two theorem identifies
the exact geometric and arithmetic images: full $\mathrm{GL}_g(2)$
for even $d\ge4$, and the full stabilizer
$2^{g-1}\!:\mathrm{GL}_{g-1}(2)$ for odd $d\ge5$.
The subsequent projective-space theorems settle every remaining
plane degree and every dimension, in all characteristics.
For arbitrary $X$, the explicit asymptotic positivity bound above is
proved using a normal-crossings union of hyperplane sections.
The later all-ambient theorem also proves every point-count residue
with an explicit degree bound. An effective finite generic-rank
calculation and exact determinant image are provided below.
A structural generic-rank formula and the complete nonabelian
monodromy group for arbitrary $X$ remain unresolved here.
